% SPDX-License-Identifier: Apache-2.0
% covers: Board
\datasheet{Board}{The Vreteno machine for the Alinx AX7A200, as one lowered unit of units.}

\dsfacts{\code{cpu/vreteno/src/board.rs}}{\code{vreteno32::board::Board}}%
{\code{//docs:vreteno}}

\subsection*{Function}

\code{Board} is everything that computes on the board, written as a
unit of units so that \code{\#[lower]} makes its netlist as one module
holding the rest. It holds the core and its tracker, a router of five,
and the peripherals behind its five ports. The data memory and the
timer each sit behind an \code{AxiPer}. Four small peripherals share
the page at \code{0x3000} behind the AXI-Lite bridge
\code{LiteBridge4}, a sixteenth of the page each: the serial port at
\code{0x3000}, the pulse width modulator at \code{0x3100}, whatever a
board hangs on the slot at \code{0x3200}, and the remote peripheral at
\code{0x3300}. The DDR3 memory sits behind an \code{AxiPer} and
holds UberDDR3's controller as a foreign module. The interrupt
controller, \code{Plic2}, sits behind a second \code{LiteBridge1}.

\code{run.rs} wires the same parts for a simulation. \code{Board} is
that wiring as a unit, so a board top has to supply only the clocks,
the resets and the pins. The top for the AX7A200 is
\code{cpu/vreteno/board/vreteno\_board.v}.

\subsection*{Parameters}

\begin{center}\footnotesize
\begin{tabular}{@{}lp{0.7\textwidth}@{}}
\toprule
Parameter & Meaning \\
\midrule
\code{DIV} & The serial port's cycles per bit, as \code{Uart<DIV>}.
The tables use 868, the value \code{board\_netlist board} writes for
115200 baud at 100~MHz. \\
\code{MICRON\_SIM} & The DDR3 controller's \code{MICRON\_SIM}: one to
shorten its power-on waits for the Micron model, zero on the board. \\
\code{BIST} & The DDR3 controller's \code{BIST\_MODE}, its self test
after calibration. \\
\bottomrule
\end{tabular}
\end{center}

\code{board\_netlist} lowers \code{Board<868, 0, 0>} for the board and
\code{Board<16, 1, 0>} for simulation. \code{//cpu/vreteno:board\_test}
runs \code{Board<4, 1, 0>}.

\subsection*{Address map}

\code{BoardRouter} states the map in its type, \code{Router5<32, 32, 4,
2, ...>}. A peripheral gets an address when the address under the mask
equals the base.

\begin{center}\footnotesize
\begin{tabular}{@{}lllp{0.35\textwidth}@{}}
\toprule
Base & Mask & Child & Behind \\
\midrule
\code{0x1000} & \code{0xffff\_f000} & \code{dmem}, \code{Dmem<2>} &
\code{pdmem}, \code{AxiPer} \\
\code{0x0200\_0000} & \code{0xffff\_0000} & \code{timer},
\code{Timer<2>} & \code{ptimer}, \code{AxiPer} \\
\code{0x3000} & \code{0xffff\_f000} & \code{uart}, \code{Uart<DIV>} &
\code{puart}, \code{LiteBridge4} at \code{0x3000} \\
& \code{0xffff\_ff00} & \code{pwm}, \code{Pwm} &
the same bridge, at \code{0x3100} \\
& \code{0xffff\_ff00} & the board's, on its own clock &
the same bridge, at \code{0x3200} \\
& \code{0xffff\_ff00} & \code{remote},
\code{Remote<REMOTE\_WAIT>} & the same bridge, at \code{0x3300} \\
\code{0x4000\_0000} & \code{0xc000\_0000} & \code{ddr3},
\code{Ddr3Per} & \code{pddr3}, \code{AxiPer} \\
\code{0x0c00\_0000} & \code{0xfc00\_0000} & \code{plic},
\code{Plic2<0>} & \code{pplic}, \code{LiteBridge1} at
\code{0x0c00\_0000} \\
\bottomrule
\end{tabular}
\end{center}

The router answers a burst that matches no range with \code{DecErr}.
The core's instruction memory is not on the bus, and the core sends no
request for an address below \code{0x1000}.

\dstables{Board}

\subsection*{Behaviour}

\begin{itemize}
\item The core's tracker is \code{AxiHost<32, 32, 4, 2, 4>}: 32-bit
addresses and words, four lanes, two-bit identifiers, four of them.
\item \code{rst} goes to the core, the timer, the serial port and the
interrupt controller. The DDR3 controller has its own active-low
\code{ddr3\_rst\_n}.
\item The interrupt controller's source 1 is the serial port's
interrupt line, and its source 2 is \code{irq}; both ask while high.
Its line is the core's external interrupt. The timer's line goes to
the core's \code{tirq}.
\item Of the core's outputs, only \code{halt} is a port. \code{instr}
and the \code{Writeback} go nowhere.
\item The join runs the timer before the core, since the core reads
the timer's line in the same step, and the serial port before the
interrupt controller before the core, for the same reason. \code{Ddr3Per} runs its controller
before its bridge for the same reason.
\item \code{calib} and the memory's pins are \code{Ddr3Per}'s,
brought out as the unit's own ports.
\item The slot at \code{0x3200} does not reach a field. It leaves as
five channel ports, \code{vaw}, \code{var} and \code{vw} out and
\code{vb} and \code{vr} back, because what sits there runs on a clock
of its own and the crossing is the board top's business. A board with
nothing there ties them off, and a program that touched \code{0x3200}
on such a board would wait for ever.
\item The remote peripheral at \code{0x3300} is the other way about:
\code{remote} and \code{link} are fields, because both run on the
core's clock. What leaves the unit is the frame stream, \code{net\_tx}
out and \code{net\_rx} in, a byte and a last bit each, which is what
the MAC speaks. \code{REMOTE\_DEV} is 1 and \code{REMOTE\_WAIT} is
two million cycles, 20~ms at 100~MHz, after which an unanswered
transaction is \code{SLVERR} on the bus.
\item \code{board\_netlist} puts a program in the core's
\code{imem} and in the data memory's \code{lane0} to \code{lane3} of
the netlist, since nothing on the machine loads them at run time.
\end{itemize}

The top, \code{vreteno\_board.v}, does the following around the
lowered \code{board}:

\begin{itemize}
\item One PLL on the board's 200~MHz differential clock makes 100~MHz
for the design, 400~MHz for \code{ddr3\_clk}, the same 400~MHz a
quarter cycle late for \code{ddr3\_clk\_90}, and 200~MHz for
\code{ref\_clk}.
\item The memory's reset follows the button and the PLL's lock. The
core's reset also waits for \code{calib}. Each passes through two
flip-flops on the design's clock.
\item It ties \code{irq} low.
\item Its four LEDs, lit when driven low, show the core halted, a
heartbeat from bit 25 of a counter, the memory calibrated, and the PLL
locked.
\end{itemize}

\subsection*{Verification}

\begin{itemize}
\item \code{//cpu/vreteno:board\_test} runs \code{Board<4, 1, 0>} in
Rust. \code{the\_greeting\_runs\_on\_the\_board} must print
\code{hello from rust} and halt within 8000 cycles.
\code{the\_memory\_test\_runs\_on\_the\_board} must print
\code{ddr3 ok} and halt within 40000 cycles.
\code{input\_comes\_by\_interrupt\_one\_byte\_each} runs a program
that takes four typed bytes through the interrupt controller, one
interrupt each, and must print \code{got ping in 4 interrupts} and
halt within 20000 cycles.
\code{the\_netlist\_holds\_the\_controller} checks that the Verilog
has the top, the core and the bridge as modules, instantiates
\code{ddr3\_wb32} without writing it, and has the \code{dq} pads.
\item \code{//cpu/vreteno/board/sim:board\_test} simulates the top, the
netlist of \code{//cpu/vreteno:board\_sim\_v}, the controller and two
Micron models under Vivado's simulator. It is manual.
\code{//docs:vreteno} reports that in that run the core says
\code{ddr3 ok} on the serial port and halts.
\item \code{//cpu/vreteno:vreteno\_board\_synth} and
\code{//cpu/vreteno:vreteno\_board\_pnr} synthesise, place and route
the design with the board's pins. Both are manual.
\code{//docs:vreteno} reports that every timing constraint is met, with
0.432~ns worst setup slack and 0.031~ns worst hold slack. It reports
8100 LUTs, 6959 flip-flops, four block RAM tiles and four DSP blocks.
\item \code{//cpu/vreteno:vreteno\_board\_prog} and
\code{//cpu/vreteno:vreteno\_board\_flash} program the FPGA and the
QSPI flash. Both are manual.
\end{itemize}

No waveform target simulates \code{Board}'s netlist against a trace.

\subsection*{Used by}

\code{//cpu/vreteno:board\_netlist}, whose genrules
\code{//cpu/vreteno:board\_v} and \code{//cpu/vreteno:board\_sim\_v}
write the netlist, and \code{vreteno\_board.v}, which instantiates it
as \code{board}. The test \code{cpu/vreteno/tests/board.rs}.

\subsection*{Limits}

\code{//docs:vreteno} states that this design, with the DDR3 memory on
the bus, has not yet been run on the board. The board run it reports
was made with an earlier design without the memory. The program is
fixed in the netlist. The board top ties \code{irq} low, so on the
board the interrupt controller's second source never asks. The core is the only host.
